% =============================================================================
% Section II — Related Work
% =============================================================================
\section{Related Work}
\label{sec:related}

\subsection{From Rule-Based Detection to Data-Driven Models}
Traditional leak detection in water systems has historically relied on
fixed thresholds, minimum night-flow analysis, and manual step-testing.
These approaches are practical and interpretable but often miss subtle
events or generate many false alarms under normal demand variation.
Machine learning methods --- support vector machines, random forests, and
k-nearest neighbours --- improved detection by learning patterns from
historical sensor data. However, shallow models depend on handcrafted lag
features and cannot natively model the multi-scale temporal dynamics that
characterise leak onset, propagation, and recovery.

\subsection{Deep Learning for Multivariate Time Series}
Deep models are now widely used for sensor anomaly detection because they
learn temporal features directly from raw sequences. One-dimensional
convolutional layers capture local change patterns such as rapid pressure
drops, while recurrent layers model longer causal context across consecutive
readings. Hybrid CNN-LSTM variants are especially effective for industrial
time-series classification where both local motifs and long-range dependencies
matter. For water distribution networks, LSTM-based models have been shown
to outperform classical baselines on real pressure sensor data when leak
events are present at the sequence level rather than the sample level.
For this reason, the present work adopts a sequential CNN-BiLSTM design
that processes per-sensor sliding windows.

\subsection{Attention and Decision Thresholding}
Self-attention mechanisms, adapted from natural language processing, allow
sequence models to produce importance-weighted context vectors that highlight
the most anomaly-relevant timesteps within each window. This property is
particularly valuable for leak detection because the hydraulic signature
of a leak is concentrated in a short segment of an otherwise normal pressure
profile. Even with a well-trained scoring model, the final operational
behaviour depends strongly on the chosen decision threshold. In leak
detection, threshold selection can be more critical than raw AUC because
operators may prioritise either high recall (zero missed leaks) or higher
precision (fewer false alarms triggering unnecessary field inspections).
This paper explicitly analyses threshold sensitivity in addition to
standard ROC and PR metrics.

\subsection{Position of This Work}
This project contributes a practical implementation and end-to-end analysis
pipeline, not only a model architecture. It integrates data cleaning,
domain-informed feature engineering, per-sensor windowing, compact sequence
modelling with temporal attention, calibration-based threshold selection,
and paper-ready visual reporting in one reproducible workflow. The model
achieves non-trivial AUC-PR above the random baseline on a severely
imbalanced 221-window test set while remaining deployable on
resource-constrained hardware at under 24\,k parameters.

